%% Auto-generated by simulations/fermi_estimate_check.py --write-tex
%% Do not edit manually. Edit the Python script instead and regenerate.
%% Frequency-band field values from Anastassiou & Koch 2011, 2015.

% --- Global parameters ---

\newcommand{\fermiLam}{0.2}           % LFP decay length (mm)
\newcommand{\fermiLamMin}{0.1}           % measured lower bound
\newcommand{\fermiLamMax}{0.3}           % measured upper bound
\newcommand{\fermiRho}{5e+04}           % neuron density (mm$^{-3}$)
\newcommand{\fermiShift}{0.4}           % ephaptic shift per mV/mm (ms)
\newcommand{\fermiGamma}{1.5}           % Lorentzian half-width of frequency spread (rad/s)
\newcommand{\fermiThreshold}{2.0}           % $K_c/\gamma$, noiseless Lorentzian

% --- Measured field range (mV/mm) across all EEG bands ---
\newcommand{\fermiFieldMin}{1}           % 
\newcommand{\fermiFieldMax}{5}           % 

% --- Band-specific field amplitudes (mV/mm) and frequencies (Hz) ---
\newcommand{\fermiEgamma}{0.5}
\newcommand{\fermiFgamma}{40}
\newcommand{\fermiEbeta}{1.0}
\newcommand{\fermiFbeta}{20}
\newcommand{\fermiEalpha}{1.5}
\newcommand{\fermiFalpha}{10}
\newcommand{\fermiEtheta}{3.0}
\newcommand{\fermiFtheta}{6}
\newcommand{\fermiEdelta}{2.0}
\newcommand{\fermiFdelta}{2}

% --- Conservative bound parameters (low end of every range) ---
\newcommand{\fermiConsE}{1.0}           % conservative field (mV/mm)
\newcommand{\fermiConsShift}{0.3}           % conservative shift per mV/mm (ms)
\newcommand{\fermiConsLam}{0.15}           % conservative decay length (mm)
\newcommand{\fermiConsGamma}{2.0}           % conservative Lorentzian width (rad/s)

%% --- Parameterized kernel functions ---
%% \fermiK{E}{f} = coupling constant K, e.g. \fermiK{\fermiEtheta}{\fermiFgamma}
%% \fermiKGamma{E}{f} = ratio K/gamma, read against K_c/gamma = 2 (noiseless Lorentzian)
%% \fermiKGammalam{E}{f}{lam} = K/gamma with explicit lam (for range statement)
%% Nothing here divides by a phase-diffusion D: the identical-frequency noisy
%% threshold K_c = 2D is a different statement and takes a different input.

\newcommand{\fermiN}{\fpeval{round(4/3*pi*\fermiLam^3*5e+04,0)}}
\newcommand{\fermiK}[2]{\fpeval{round(4/3*pi*\fermiLam^3*5e+04 * (#1*0.4/1000) * #2, 0)}}
\newcommand{\fermiKGamma}[2]{\fpeval{round(4/3*pi*\fermiLam^3*5e+04 * (#1*0.4/1000) * #2 / \fermiGamma, 1)}}
\newcommand{\fermiKGammalam}[3]{\fpeval{round(4/3*pi*#3^3*5e+04 * (#1*0.4/1000) * #2 / \fermiGamma, 1)}}

% --- Pre-computed values used directly in prose ---
\newcommand{\fermiPhAdv}{\fpeval{round(3.0*0.4/1000*6,3)}}
\newcommand{\fermiKGammaCons}{\fpeval{round(4/3*pi*\fermiConsLam^3*5e+04 * (1.0*0.4/1000) * 40 / \fermiConsGamma, 1)}}
